% ============================================================
%  v4.2.0-r01 development patch
%  OP1/OP4 Joint HC-Dixon Reduction
%  Scope: source-only DEV-PATCH; no release claim, no PDF freeze.
% ============================================================

\subsection{OP1/OP4 joint HC--Dixon reduction}
\label{ssec:op1-op4-joint-hcdixon-reduction}

\begin{remark}[Scope of the OP1/OP4 reduction]
\label{rem:op1-op4-reduction-scope}
The purpose of this subsection is not to claim a global unconditional
solution of OP1 or OP4.  It records a local closure architecture inside
the declared OP5-inherited HC--Halo/Dixon grammar.  The construction
reduces OP1 to a continuum-compression determinant and reduces the
leading CKM/PMNS part of OP4 to a finite Fano/Dixon selector normal form
over the same compression factor.  The remaining proof obligations are
kept explicit as residual registers.
\end{remark}

\begin{definition}[HS(0)-anchored Hilbert-shift ladder]
\label{def:hs0-anchored-ladder}
The unshifted layer
\[
  \mathrm{HS}(0)
\]
is the raw kernel aim, i.e. the direct unshifted continuum/collapse
target before HC capture.  It carries unit transport weight.  The active
pre-Dixon compression layers are
\[
  \mathrm{HS}(1),\ldots,\mathrm{HS}(5),
\]
with the following structural readings:
\[
\begin{array}{c|l}
\mathrm{HS}(1) & \text{HC collar / no-through-collapse boundary},\\
\mathrm{HS}(2) & \text{normal orbit / stable collar},\\
\mathrm{HS}(3) & \text{Fano fibre / octonionic incidence},\\
\mathrm{HS}(4) & M_{16}/P_{C_7}\text{ meta-septum leakage gate},\\
\mathrm{HS}(5) & \text{two-sided front/back carrier}.
\end{array}
\]
The next layer
\[
  \mathrm{HS}(6)
\]
is the Dixon addressing carrier
\[
  \mathbb C\otimes\mathbb H\otimes\mathbb O,
\]
and is therefore not counted as a further pre-Dixon compression layer.
\end{definition}

\begin{lemma}[HS(0)-anchored pre-Dixon capacity]
\label{lem:hs0-predixon-capacity}
Under the HS(0)-anchored ladder, the active pre-Dixon compression
capacity is
\[
  \Gamma_{\mathrm{preDix}}
  =
  \prod_{k=1}^{5}\gL
  =
  \gL^5.
\]
The layer \(\mathrm{HS}(0)\) has weight \(1\), while \(\mathrm{HS}(6)\)
is the Dixon readout/addressing carrier.
\end{lemma}

\begin{definition}[Continuum-limit compression map]
\label{def:cl-compression-map}
Let
\[
  \Xi_{\mathrm{CL}}^{\mathrm{HC/Dix}}
  :
  \mathcal R_{\mathrm{raw}}^{\mathrm{CL}}
  \longrightarrow
  \mathcal R_{\mathrm{adm}}^{\mathrm{HC/Dix}}
\]
be the raw-to-admissible continuum-limit readout map inside the
HC/Dixon grammar.  Define the continuum-limit compression factor by
\[
  \delta_{\mathrm{CL}}
  :=
  \det_{\mathrm{norm}}
  \left(
    D\Xi_{\mathrm{CL}}^{\mathrm{HC/Dix}}
  \right).
\]
\end{definition}

\begin{definition}[OP1 compression load]
\label{def:op1-compression-load}
The OP1 compression load is the normalized first spectral raw load
\[
  \Omega_{\mathrm{CL}}
  :=
  \frac{\gamma_1}{4\pi\gL^5},
\]
where \(\gamma_1\) is the first non-trivial Riemann-zero height,
\(4\pi\) is the spherical collar normalization, and \(\gL^5\) is the
five-stage pre-Dixon capacity of Lemma~\ref{lem:hs0-predixon-capacity}.
\end{definition}

\begin{lemma}[Exponential HC/Dixon determinant]
\label{lem:exponential-hcdixon-determinant}
Let
\[
  \Xi_t: \mathcal R_{\mathrm{raw}}\to \mathcal R_{\mathrm{adm}},
  \qquad 0\le t\le 1,
\]
be a route-faithful HC/Dixon readout homotopy.  Suppose the local
admissibility-loss generator \(A(t)\) satisfies
\[
  \frac{d}{dt}
  \log
  \det_{\mathrm{norm}}D\Xi_t
  =
  -\operatorname{tr}A(t).
\]
If
\[
  \Omega
  =
  \int_0^1 \operatorname{tr}A(t)\,dt,
\]
then
\[
  \det_{\mathrm{norm}}D\Xi_1
  =
  e^{-\Omega}.
\]
\end{lemma}

\begin{corollary}[OP1 compression factor]
\label{cor:op1-compression-factor}
Assuming the OP1 load identity of Definition~\ref{def:op1-compression-load}
and the exponential determinant lemma,
\[
  \delta_{\mathrm{CL}}
  =
  \exp\left(
    -\frac{\gamma_1}{4\pi\gL^5}
  \right).
\]
Consequently,
\[
  \bsym
  =
  \frac{1}{\sqrt3}\delta_{\mathrm{CL}}
  =
  \frac{1}{\sqrt3}
  \exp\left(
    -\frac{\gamma_1}{4\pi\gL^5}
  \right).
\]
\end{corollary}

\begin{definition}[CL-to-OP5 load factorization]
\label{def:cl-to-op5-factorization}
Let
\[
  \Xi_{\mathrm{OP5}}^{\mathrm{Dix}}
  :
  \mathcal R_{\mathrm{nearK}}^{\mathrm{OP5}}
  \longrightarrow
  \{\mathsf{adm},\mathsf{block},\mathsf{abs},\mathsf{route}\}
\]
be the residual-free OP5 HC/Dixon route classifier.  A CL-to-OP5
factorization is a representation
\[
  \Xi_{\mathrm{CL}}^{\mathrm{Dix}}
  =
  \mathcal C_{\mathrm{load}}
  \circ
  \Xi_{\mathrm{OP5}}^{\mathrm{Dix}}
  \big|_{\mathcal R_{\mathrm{raw}}^{\mathrm{CL}}},
\]
where
\[
  \mathcal R_{\mathrm{raw}}^{\mathrm{CL}}
  \subseteq
  \mathcal R_{\mathrm{nearK}}^{\mathrm{OP5}},
\]
and
\[
  \mathcal C_{\mathrm{load}}
  :
  \{\mathsf{adm},\mathsf{block},\mathsf{abs},\mathsf{route}\}
  \to
  \mathbb R_{\ge0}
\]
assigns compression load to classified OP5 route types.
\end{definition}

\begin{lemma}[CL residual inheritance from OP5]
\label{lem:cl-residual-inheritance-op5}
Assume the OP5 classifier has no residual class,
\[
  \mathcal I_{\mathrm{res}}^{\mathrm{OP5/Dix}}=\varnothing,
\]
and assume the CL-to-OP5 factorization of
Definition~\ref{def:cl-to-op5-factorization}.  Then the continuum-limit
residual register is empty:
\[
  \mathcal R_{\mathrm{res}}^{\mathrm{CL/Dix}}=\varnothing.
\]
\end{lemma}

\begin{definition}[Fano/Dixon mediator table]
\label{def:fano-dixon-mediator-table}
Let the quark generation carriers be represented by
\[
  \mathrm{Gen}_1\sim e_4,
  \qquad
  \mathrm{Gen}_2\sim e_5,
  \qquad
  \mathrm{Gen}_3\sim e_6.
\]
The three Fano products are
\[
  e_4e_5=e_7,\qquad
  e_5e_6=e_1,\qquad
  e_6e_4=-e_3.
\]
They are read as mediator types
\[
\begin{array}{c|c|c}
\text{channel} & \text{mediator} & \text{type}\\
\hline
12 & e_7 & \text{HS(4)-gated late Fano edge},\\
23 & e_1 & \text{quaternionic-over-triolic lift},\\
13 & -e_3 & \text{orientation-inverted frame mediator}.
\end{array}
\]
\end{definition}

\begin{definition}[Finite Fano/Dixon selectors]
\label{def:fano-dixon-selectors}
The finite selector weights assigned to the mediator types are
\[
  \chi(e_7)=\frac{15}{16},
  \qquad
  \chi(e_1)=\frac43,
  \qquad
  \chi(-e_3)=\delta_{\mathrm{CL}}.
\]
The last selector means that the \(13\)-quark channel carries one
additional frame-compression factor besides the universal
\(\delta_{\mathrm{CL}}\).
\end{definition}

\begin{lemma}[Quark-sector selector normal form]
\label{lem:quark-selector-normal-form}
Let
\[
  \gamma_s=\frac{\ln2}{\pi\sqrt3}.
\]
Assign Fano path lengths
\[
  d_{12}=1,\qquad d_{23}=2,\qquad d_{13}=3.
\]
Then the leading quark-sector readouts are
\[
  s_{ij}^{q}
  =
  2\gamma_s^{d_{ij}}
  \chi(\mathfrak m_{ij})
  \delta_{\mathrm{CL}},
\]
where \(\mathfrak m_{ij}\) is the mediator of the \(ij\)-channel.
Equivalently,
\[
  s_{12}^{q}
  =
  2\cdot\frac{15}{16}\gamma_s\delta_{\mathrm{CL}},
\]
\[
  s_{23}^{q}
  =
  2\cdot\frac43\gamma_s^2\delta_{\mathrm{CL}},
\]
\[
  s_{13}^{q}
  =
  2\gamma_s^3\delta_{\mathrm{CL}}^2.
\]
\end{lemma}

\begin{theorem}[Leading CKM Fano/Dixon frame-readout candidate]
\label{thm:leading-ckm-fano-dixon-candidate}
Assume the quark-sector selector normal form of
Lemma~\ref{lem:quark-selector-normal-form}.  Define
\[
  \Omega_{\mathrm{CL}}=-\log\delta_{\mathrm{CL}}
\]
and
\[
  \phi_q
  =
  \frac{\pi}{3}+2\Omega_{\mathrm{CL}}.
\]
Then
\[
  V_{\mathrm{CKM}}^{\mathrm{FD}}
  =
  R_{23}(s_{23}^{q})
  R_{13}(s_{13}^{q},\phi_q)
  R_{12}(s_{12}^{q})
\]
is unitary by construction and gives a leading CKM readout candidate
inside the Fano/Dixon/frame grammar.
\end{theorem}

\begin{remark}[CKM status]
\label{rem:ckm-status-fd}
The theorem is a local readout candidate, not a global derivation of
the CKM matrix outside the declared grammar.  Its proof obligations are
the HS(4)-gate value \(15/16\), the quaternionic-over-triolic lift
\(4/3\), the orientation-inverted \(13\)-channel compression, the OP1
compression factor, and the frame-extraction map to standard matter-frame
fit values.
\end{remark}

\begin{definition}[Triolic PMNS base angles]
\label{def:triolic-pmns-base-angles}
Define
\[
  \alpha_T=\arcsin\left(\frac1{\sqrt3}\right),
  \qquad
  \beta_T=\arctan(\sqrt2).
\]
\end{definition}

\begin{lemma}[PMNS triolic/Fano-leakage readout candidate]
\label{lem:pmns-triolic-fano-leakage}
The leading PMNS angle candidates are
\[
  \theta_{12}^{\ell}
  =
  \alpha_T\delta_{\mathrm{CL}},
\]
\[
  \theta_{23}^{\ell}
  =
  \beta_T\delta_{\mathrm{CL}}^2,
\]
and
\[
  \theta_{13}^{\ell}
  =
  \arcsin\left(\frac76\gamma_s\right).
\]
The factor
\[
  \frac76
  =
  \frac{|\operatorname{Im}\mathbb O|}
       {|\operatorname{OrEdges}(K_3)|}
\]
is read as the minimal Fano-over-oriented-edge leakage selector.
\end{lemma}

\begin{theorem}[Leading PMNS triolic/Fano readout candidate]
\label{thm:leading-pmns-fd-candidate}
Let
\[
  U_{\mathrm{PMNS}}^{\mathrm{FD}}
  =
  R_{23}(\theta_{23}^{\ell})
  R_{13}(\theta_{13}^{\ell},\phi_{\ell})
  R_{12}(\theta_{12}^{\ell}).
\]
Then \(U_{\mathrm{PMNS}}^{\mathrm{FD}}\) is a leading PMNS readout
candidate in which the solar and atmospheric angles arise from
triolic fixed-point geometry and the reactor angle arises from a
Fano-over-edge leakage selector.  The leptonic CP phase
\(\phi_{\ell}\) remains an ordering-/matter-/propagation-sensitive branch.
\end{theorem}

\begin{definition}[Frame-extraction residual]
\label{def:frame-extraction-residual}
Let \(O_{\mathrm{FD}}\) be a HC/Dixon/Fano/triolic readout and
\(O_{\mathrm{std}}\) the corresponding standard matter-frame
global-fit observable.  A frame-extraction map is a map
\[
  \mathcal E_{\mathrm{frame}}
\]
such that
\[
  O_{\mathrm{std}}
  =
  \mathcal E_{\mathrm{frame}}(O_{\mathrm{FD}}).
\]
The frame residual is
\[
  \rho_{\mathrm{frame}}
  =
  O_{\mathrm{std}}-O_{\mathrm{FD}}.
\]
\end{definition}

\begin{remark}[Frame-sensitive residuals]
\label{rem:frame-sensitive-residuals}
The residual \(\rho_{\mathrm{frame}}\) is expected to split into
ordinary frame compression, orientation-sensitive compression,
phase-load correction, fit extraction, matter-effect, ordering, and
octant contributions.  Thus percent-level residuals against standard
global-fit values are not automatically contradictions; they must be
localized in the appropriate frame/readout register.
\end{remark}

\begin{theorem}[Conditional OP1/OP4 local closure candidate]
\label{thm:conditional-op1-op4-local-closure}
Assume:
\begin{enumerate}
  \item the OP5 HC/Dixon route classifier is residual-free;
  \item the CL-to-OP5 factorization holds;
  \item the HS(0)-anchored pre-Dixon capacity is \(\gL^5\);
  \item the OP1 load identity is
  \[
    \Omega_{\mathrm{CL}}
    =
    \frac{\gamma_1}{4\pi\gL^5};
  \]
  \item the exponential determinant lemma applies;
  \item the finite Fano/Dixon selector identities
  \[
    15/16,
    \qquad
    4/3,
    \qquad
    7/6
  \]
  hold in the indicated channels.
\end{enumerate}
Then OP1 locally closes inside the OP5-inherited HC/Dixon
continuum-compression grammar, and the leading CKM/PMNS sector of OP4
is reduced to finite Fano/Dixon/triolic readout data over the same
compression factor \(\delta_{\mathrm{CL}}\).
\end{theorem}

\begin{remark}[Non-global status of the joint closure candidate]
\label{rem:op1-op4-nonglobal-status}
The conclusion of Theorem~\ref{thm:conditional-op1-op4-local-closure}
is local to the declared grammar.  It does not assert a global
unconditional derivation of CKM or PMNS data, nor a global OP1/OP4
closure outside the HC/Dixon/Fano/triolic frame-readout setting.
The remaining proof obligations are finite and explicit.
\end{remark}

\begin{definition}[Joint OP1/OP4 residual ledger]
\label{def:op1-op4-residual-ledger}
Define the joint residual vector
\[
  \rho_{\mathrm{OP1/4}}
  =
  ( 
    \rho_{\mathrm{fact}},
    \rho_{\mathrm{CL}},
    \rho_{15/16},
    \rho_{4/3},
    \rho_{7/6},
    \rho_{\phi_q},
    \rho_{\phi_{\ell}},
    \rho_{\mathrm{frame}}
  ).
\]
Here
\[
  \rho_{\mathrm{fact}}
  =
  \mathcal R_{\mathrm{res}}^{\mathrm{CL/Dix}},
\]
\[
  \rho_{\mathrm{CL}}
  =
  \delta_{\mathrm{CL}}
  -
  \exp\left(
    -\frac{\gamma_1}{4\pi\gL^5}
  \right),
\]
\[
  \rho_{15/16}
  =
  \chi(e_7)-\frac{15}{16},
  \qquad
  \rho_{4/3}
  =
  \chi(e_1)-\frac43,
  \qquad
  \rho_{7/6}
  =
  \chi_{\ell13}-\frac76,
\]
and
\[
  \rho_{\phi_q}
  =
  \phi_q-\\left(\frac{\pi}{3}+2\Omega_{\mathrm{CL}}\right),
\]
while \(\rho_{\phi_{\ell}}\) and \(\rho_{\mathrm{frame}}\) record the
leptonic CP branch and the frame-extraction residuals.
\end{definition}

\begin{remark}[Closure condition]
\label{rem:op1-op4-ledger-closure}
If
\[
  \rho_{\mathrm{OP1/4}}=0,
\]
then OP1 and the leading CKM/PMNS part of OP4 locally close inside the
declared OP5-inherited HC/Dixon/Fano/triolic readout grammar.
\end{remark}

\begin{remark}[Numerical audit status]
\label{rem:op1-op4-numerical-audit-status}
The numerical agreement of the leading CKM and PMNS readouts is read as
a consistency audit of the finite register architecture, not as a
substitute for the missing selector and frame-extraction proofs.  The
same compression factor \(\delta_{\mathrm{CL}}\) appears in the OP1
readout, in the CKM hierarchy, in the CKM CP-phase shift through
\(\Omega_{\mathrm{CL}}\), and in the PMNS triolic fixed-point readouts.
This repeated appearance is evidence for a common HC/Dixon
frame-compression layer, but it is not an unconditional theorem.
\end{remark}
